% --- Archon physics formalization source begin ---
% archon:physics
% archon:covers PhyXMiniProblems/problem_phyx_mini_0783.lean
% archon:source-report reports/phyx_mini/problem_phyx_mini_0783.source.json

\chapter{Physics problem phyx\_mini\_0783}
\label{ch:PhyXMiniProblems_problem_phyx_mini_0783}

\paragraph{Problem source.}
\#\# Physical scenario

figure illustrates an Atwood's machine. if there is no slipping between the cord and the surface of the wheel. Let the masses of blocks \textbackslash{}(A\textbackslash{}) and \textbackslash{}(B\textbackslash{}) be \textbackslash{}(4.00\textbackslash{}text\{ kg\}\textbackslash{}) and \textbackslash{}(2.00\textbackslash{}text\{ kg\}\textbackslash{}), respectively, the moment of inertia of the wheel about its axis be \textbackslash{}(0.220\textbackslash{}text\{ kg\}\textbackslash{}cdot\textbackslash{}text\{m\}\textasciicircum{}2\textbackslash{}), and the radius of the wheel be \textbackslash{}(0.120\textbackslash{}text\{ m\}\textbackslash{}).

\#\# Question

Find the angular acceleration of the wheel \textbackslash{}(C\textbackslash{}).

\#\# Answer choices

- A: A: \textbackslash{}(12.14\textbackslash{}text\{ rad/s\}\textasciicircum{}2\textbackslash{})

- B: B: \textbackslash{}(7.68\textbackslash{}text\{ rad/s\}\textasciicircum{}2\textbackslash{})

- C: C: \textbackslash{}(5.28\textbackslash{}text\{ rad/s\}\textasciicircum{}2\textbackslash{})

- D: D: \textbackslash{}(6.68\textbackslash{}text\{ rad/s\}\textasciicircum{}2\textbackslash{})

\#\# Figure caption (auxiliary; use the image as primary evidence)

The image shows a mechanical setup with a pulley system. There is a single pulley labeled "C" suspended from above. Two masses, labeled "A" and "B," are hanging from the pulley. Mass "A" is larger than mass "B" and both are rectangular in shape. Each mass is attached to a string that runs over the pulley, creating a simple pulley arrangement for lifting or balancing the masses. The strings appear taut, indicating tension in the system.

\#\# Dataset metadata

Rotational Motion; Physical Model Grounding Reasoning, Multi-Formula Reasoning

\paragraph{Recorded answer/context.}
B: B: \textbackslash{}(7.68\textbackslash{}text\{ rad/s\}\textasciicircum{}2\textbackslash{})

\paragraph{Figure/image path.}
/root/proposal\_for\_physic/science-mango/phyx\_mini\_run/phyx\_data/test\_image/783.png

\paragraph{Formalization target.}
create a compiling Lean file with sorry bodies at `PhyXMiniProblems/problem\_phyx\_mini\_0783.lean`. The Lean declarations must preserve the physical quantities, dimensions or dimensional roles, figure labels, governing-law hypotheses, and final relation expressed by this problem.
Use Mathlib/Physlib names found through LeanExplore where available. If a physics API is missing, introduce faithful local abstractions rather than scalar placeholder aliases.

\begin{theorem}[Physics formalization target]
\label{thm:physics:phyx_mini_0783:target}
\lean{PhyXMiniProblems.ProblemPhyXMini0783.atwoodWheelAngularAcceleration}
\uses{def:physics:phyx-mini-0783:phyxminiproblems-problemphyxmini0783-angularaccelerationinradianspersecondsquared, def:physics:phyx-mini-0783:phyxminiproblems-problemphyxmini0783-atwoodwheelsetup, def:physics:phyx-mini-0783:phyxminiproblems-problemphyxmini0783-matchesatwoodscenario, def:physics:phyx-mini-0783:phyxminiproblems-problemphyxmini0783-matchesproblemreadouts, def:physics:phyx-mini-0783:phyxminiproblems-problemphyxmini0783-matchessuppliedatwoodfigure, def:physics:phyx-mini-0783:phyxminiproblems-problemphyxmini0783-matchesterrestrialgravitycalibration, def:physics:phyx-mini-0783:phyxminiproblems-problemphyxmini0783-hasphysicalatwoodparameters, def:physics:phyx-mini-0783:phyxminiproblems-problemphyxmini0783-satisfiesatwoodwheeldynamics, def:physics:phyx-mini-0783:phyxminiproblems-problemphyxmini0783-matchesdisplayedanswer}
The assigned autoformalize agent should translate this physics problem into Lean declarations in the covered file, with theorem and lemma proof bodies written as `by sorry`.
\end{theorem}
\begin{proof}
This is an autoformalization task, not a proof task. Produce statements that can later be proved by the physics prover without weakening the physical model.
\end{proof}

% --- Archon declaration topology begin ---
\section{Declaration topology}
\begin{definition}[acceleration Dimension]
  \label{def:physics:phyx-mini-0783:phyxminiproblems-problemphyxmini0783-accelerationdimension}
  \lean{PhyXMiniProblems.ProblemPhyXMini0783.accelerationDimension}
  Linear acceleration has physical dimension L T⁻²
\end{definition}
\begin{proof}
  This is the declared model definition of acceleration Dimension.
\end{proof}
\begin{definition}[force Dimension]
  \label{def:physics:phyx-mini-0783:phyxminiproblems-problemphyxmini0783-forcedimension}
  \lean{PhyXMiniProblems.ProblemPhyXMini0783.forceDimension}
  Force has physical dimension M L T⁻²
\end{definition}
\begin{proof}
  This is the declared model definition of force Dimension.
\end{proof}
\begin{definition}[moment Of Inertia Dimension]
  \label{def:physics:phyx-mini-0783:phyxminiproblems-problemphyxmini0783-momentofinertiadimension}
  \lean{PhyXMiniProblems.ProblemPhyXMini0783.momentOfInertiaDimension}
  A scalar moment of inertia has physical dimension M L²
\end{definition}
\begin{proof}
  This is the declared model definition of moment Of Inertia Dimension.
\end{proof}
\begin{definition}[angular Acceleration Dimension]
  \label{def:physics:phyx-mini-0783:phyxminiproblems-problemphyxmini0783-angularaccelerationdimension}
  \lean{PhyXMiniProblems.ProblemPhyXMini0783.angularAccelerationDimension}
  Radians are dimensionless, so angular acceleration has dimension T⁻²
\end{definition}
\begin{proof}
  This is the declared model definition of angular Acceleration Dimension.
\end{proof}
\begin{definition}[Mass Quantity]
  \label{def:physics:phyx-mini-0783:phyxminiproblems-problemphyxmini0783-massquantity}
  \lean{PhyXMiniProblems.ProblemPhyXMini0783.MassQuantity}
  A nonnegative, unit-independent physical mass
\end{definition}
\begin{proof}
  This is the declared model definition of Mass Quantity.
\end{proof}
\begin{definition}[Length Quantity]
  \label{def:physics:phyx-mini-0783:phyxminiproblems-problemphyxmini0783-lengthquantity}
  \lean{PhyXMiniProblems.ProblemPhyXMini0783.LengthQuantity}
  A nonnegative, unit-independent physical length
\end{definition}
\begin{proof}
  This is the declared model definition of Length Quantity.
\end{proof}
\begin{definition}[Acceleration Magnitude Quantity]
  \label{def:physics:phyx-mini-0783:phyxminiproblems-problemphyxmini0783-accelerationmagnitudequantity}
  \lean{PhyXMiniProblems.ProblemPhyXMini0783.AccelerationMagnitudeQuantity}
  \uses{def:physics:phyx-mini-0783:phyxminiproblems-problemphyxmini0783-accelerationdimension}
  A nonnegative linear-acceleration magnitude
\end{definition}
\begin{proof}
  This is the declared model definition of Acceleration Magnitude Quantity.
\end{proof}
\begin{definition}[Force Magnitude Quantity]
  \label{def:physics:phyx-mini-0783:phyxminiproblems-problemphyxmini0783-forcemagnitudequantity}
  \lean{PhyXMiniProblems.ProblemPhyXMini0783.ForceMagnitudeQuantity}
  \uses{def:physics:phyx-mini-0783:phyxminiproblems-problemphyxmini0783-forcedimension}
  A nonnegative force magnitude, used for the two cord tensions
\end{definition}
\begin{proof}
  This is the declared model definition of Force Magnitude Quantity.
\end{proof}
\begin{definition}[Moment Of Inertia Quantity]
  \label{def:physics:phyx-mini-0783:phyxminiproblems-problemphyxmini0783-momentofinertiaquantity}
  \lean{PhyXMiniProblems.ProblemPhyXMini0783.MomentOfInertiaQuantity}
  \uses{def:physics:phyx-mini-0783:phyxminiproblems-problemphyxmini0783-momentofinertiadimension}
  A nonnegative moment of inertia about the wheel axle
\end{definition}
\begin{proof}
  This is the declared model definition of Moment Of Inertia Quantity.
\end{proof}
\begin{definition}[Angular Acceleration Magnitude Quantity]
  \label{def:physics:phyx-mini-0783:phyxminiproblems-problemphyxmini0783-angularaccelerationmagnitudequantity}
  \lean{PhyXMiniProblems.ProblemPhyXMini0783.AngularAccelerationMagnitudeQuantity}
  \uses{def:physics:phyx-mini-0783:phyxminiproblems-problemphyxmini0783-angularaccelerationdimension}
  A nonnegative angular-acceleration magnitude
\end{definition}
\begin{proof}
  This is the declared model definition of Angular Acceleration Magnitude Quantity.
\end{proof}
\begin{definition}[mass Readout]
  \label{def:physics:phyx-mini-0783:phyxminiproblems-problemphyxmini0783-massreadout}
  \lean{PhyXMiniProblems.ProblemPhyXMini0783.massReadout}
  \uses{def:physics:phyx-mini-0783:phyxminiproblems-problemphyxmini0783-massquantity}
  Read a mass in the coherent mass unit selected by units
\end{definition}
\begin{proof}
  This is the declared model definition of mass Readout.
\end{proof}
\begin{definition}[length Readout]
  \label{def:physics:phyx-mini-0783:phyxminiproblems-problemphyxmini0783-lengthreadout}
  \lean{PhyXMiniProblems.ProblemPhyXMini0783.lengthReadout}
  \uses{def:physics:phyx-mini-0783:phyxminiproblems-problemphyxmini0783-lengthquantity}
  Read a length in the coherent length unit selected by units
\end{definition}
\begin{proof}
  This is the declared model definition of length Readout.
\end{proof}
\begin{definition}[acceleration Readout]
  \label{def:physics:phyx-mini-0783:phyxminiproblems-problemphyxmini0783-accelerationreadout}
  \lean{PhyXMiniProblems.ProblemPhyXMini0783.accelerationReadout}
  \uses{def:physics:phyx-mini-0783:phyxminiproblems-problemphyxmini0783-accelerationmagnitudequantity}
  Read linear acceleration in the coherent units selected by units
\end{definition}
\begin{proof}
  This is the declared model definition of acceleration Readout.
\end{proof}
\begin{definition}[force Readout]
  \label{def:physics:phyx-mini-0783:phyxminiproblems-problemphyxmini0783-forcereadout}
  \lean{PhyXMiniProblems.ProblemPhyXMini0783.forceReadout}
  \uses{def:physics:phyx-mini-0783:phyxminiproblems-problemphyxmini0783-forcemagnitudequantity}
  Read a force magnitude in the coherent units selected by units
\end{definition}
\begin{proof}
  This is the declared model definition of force Readout.
\end{proof}
\begin{definition}[moment Of Inertia Readout]
  \label{def:physics:phyx-mini-0783:phyxminiproblems-problemphyxmini0783-momentofinertiareadout}
  \lean{PhyXMiniProblems.ProblemPhyXMini0783.momentOfInertiaReadout}
  \uses{def:physics:phyx-mini-0783:phyxminiproblems-problemphyxmini0783-momentofinertiaquantity}
  Read moment of inertia in coherent mass and length units
\end{definition}
\begin{proof}
  This is the declared model definition of moment Of Inertia Readout.
\end{proof}
\begin{definition}[angular Acceleration Readout]
  \label{def:physics:phyx-mini-0783:phyxminiproblems-problemphyxmini0783-angularaccelerationreadout}
  \lean{PhyXMiniProblems.ProblemPhyXMini0783.angularAccelerationReadout}
  \uses{def:physics:phyx-mini-0783:phyxminiproblems-problemphyxmini0783-angularaccelerationmagnitudequantity}
  Read angular acceleration in inverse-square selected-time units
\end{definition}
\begin{proof}
  This is the declared model definition of angular Acceleration Readout.
\end{proof}
\begin{definition}[mass In Kilograms]
  \label{def:physics:phyx-mini-0783:phyxminiproblems-problemphyxmini0783-massinkilograms}
  \lean{PhyXMiniProblems.ProblemPhyXMini0783.massInKilograms}
  \uses{def:physics:phyx-mini-0783:phyxminiproblems-problemphyxmini0783-massquantity, def:physics:phyx-mini-0783:phyxminiproblems-problemphyxmini0783-massreadout}
  SI kilogram readout of a mass
\end{definition}
\begin{proof}
  This is the declared model definition of mass In Kilograms.
\end{proof}
\begin{definition}[length In Meters]
  \label{def:physics:phyx-mini-0783:phyxminiproblems-problemphyxmini0783-lengthinmeters}
  \lean{PhyXMiniProblems.ProblemPhyXMini0783.lengthInMeters}
  \uses{def:physics:phyx-mini-0783:phyxminiproblems-problemphyxmini0783-lengthquantity, def:physics:phyx-mini-0783:phyxminiproblems-problemphyxmini0783-lengthreadout}
  SI metre readout of a length
\end{definition}
\begin{proof}
  This is the declared model definition of length In Meters.
\end{proof}
\begin{definition}[acceleration In Meters Per Second Squared]
  \label{def:physics:phyx-mini-0783:phyxminiproblems-problemphyxmini0783-accelerationinmeterspersecondsquared}
  \lean{PhyXMiniProblems.ProblemPhyXMini0783.accelerationInMetersPerSecondSquared}
  \uses{def:physics:phyx-mini-0783:phyxminiproblems-problemphyxmini0783-accelerationmagnitudequantity, def:physics:phyx-mini-0783:phyxminiproblems-problemphyxmini0783-accelerationreadout}
  SI metre-per-second-squared readout of linear acceleration
\end{definition}
\begin{proof}
  This is the declared model definition of acceleration In Meters Per Second Squared.
\end{proof}
\begin{definition}[force In Newtons]
  \label{def:physics:phyx-mini-0783:phyxminiproblems-problemphyxmini0783-forceinnewtons}
  \lean{PhyXMiniProblems.ProblemPhyXMini0783.forceInNewtons}
  \uses{def:physics:phyx-mini-0783:phyxminiproblems-problemphyxmini0783-forcemagnitudequantity, def:physics:phyx-mini-0783:phyxminiproblems-problemphyxmini0783-forcereadout}
  SI newton readout of a force magnitude
\end{definition}
\begin{proof}
  This is the declared model definition of force In Newtons.
\end{proof}
\begin{definition}[moment Of Inertia In Kilogram Meters Squared]
  \label{def:physics:phyx-mini-0783:phyxminiproblems-problemphyxmini0783-momentofinertiainkilogrammeterssquared}
  \lean{PhyXMiniProblems.ProblemPhyXMini0783.momentOfInertiaInKilogramMetersSquared}
  \uses{def:physics:phyx-mini-0783:phyxminiproblems-problemphyxmini0783-momentofinertiaquantity, def:physics:phyx-mini-0783:phyxminiproblems-problemphyxmini0783-momentofinertiareadout}
  SI kilogram-square-metre readout of moment of inertia
\end{definition}
\begin{proof}
  This is the declared model definition of moment Of Inertia In Kilogram Meters Squared.
\end{proof}
\begin{definition}[angular Acceleration In Radians Per Second Squared]
  \label{def:physics:phyx-mini-0783:phyxminiproblems-problemphyxmini0783-angularaccelerationinradianspersecondsquared}
  \lean{PhyXMiniProblems.ProblemPhyXMini0783.angularAccelerationInRadiansPerSecondSquared}
  \uses{def:physics:phyx-mini-0783:phyxminiproblems-problemphyxmini0783-angularaccelerationmagnitudequantity, def:physics:phyx-mini-0783:phyxminiproblems-problemphyxmini0783-angularaccelerationreadout}
  SI readout of angular acceleration, numerically in radians per second squared
\end{definition}
\begin{proof}
  This is the declared model definition of angular Acceleration In Radians Per Second Squared.
\end{proof}
\begin{definition}[Block]
  \label{def:physics:phyx-mini-0783:phyxminiproblems-problemphyxmini0783-block}
  \lean{PhyXMiniProblems.ProblemPhyXMini0783.Block}
  The two hanging blocks carrying the figure labels A and B
\end{definition}
\begin{proof}
  This is the declared model definition of Block.
\end{proof}
\begin{definition}[Figure Label]
  \label{def:physics:phyx-mini-0783:phyxminiproblems-problemphyxmini0783-figurelabel}
  \lean{PhyXMiniProblems.ProblemPhyXMini0783.FigureLabel}
  Labels printed in the supplied raster
\end{definition}
\begin{proof}
  This is the declared model definition of Figure Label.
\end{proof}
\begin{definition}[Horizontal Side]
  \label{def:physics:phyx-mini-0783:phyxminiproblems-problemphyxmini0783-horizontalside}
  \lean{PhyXMiniProblems.ProblemPhyXMini0783.HorizontalSide}
  Horizontal positions relative to the wheel in the front-view diagram
\end{definition}
\begin{proof}
  This is the declared model definition of Horizontal Side.
\end{proof}
\begin{definition}[Vertical Direction]
  \label{def:physics:phyx-mini-0783:phyxminiproblems-problemphyxmini0783-verticaldirection}
  \lean{PhyXMiniProblems.ProblemPhyXMini0783.VerticalDirection}
  Positive directions used for scalar block equations
\end{definition}
\begin{proof}
  This is the declared model definition of Vertical Direction.
\end{proof}
\begin{definition}[Rotation Direction]
  \label{def:physics:phyx-mini-0783:phyxminiproblems-problemphyxmini0783-rotationdirection}
  \lean{PhyXMiniProblems.ProblemPhyXMini0783.RotationDirection}
  Positive sense used for the wheel's angular acceleration
\end{definition}
\begin{proof}
  This is the declared model definition of Rotation Direction.
\end{proof}
\begin{definition}[Cord Model]
  \label{def:physics:phyx-mini-0783:phyxminiproblems-problemphyxmini0783-cordmodel}
  \lean{PhyXMiniProblems.ProblemPhyXMini0783.CordModel}
  Idealization of the single cord joining the two blocks
\end{definition}
\begin{proof}
  This is the declared model definition of Cord Model.
\end{proof}
\begin{definition}[Cord Wheel Contact Model]
  \label{def:physics:phyx-mini-0783:phyxminiproblems-problemphyxmini0783-cordwheelcontactmodel}
  \lean{PhyXMiniProblems.ProblemPhyXMini0783.CordWheelContactModel}
  Contact model between the cord and the wheel surface
\end{definition}
\begin{proof}
  This is the declared model definition of Cord Wheel Contact Model.
\end{proof}
\begin{definition}[Wheel Axle Model]
  \label{def:physics:phyx-mini-0783:phyxminiproblems-problemphyxmini0783-wheelaxlemodel}
  \lean{PhyXMiniProblems.ProblemPhyXMini0783.WheelAxleModel}
  Support and bearing model for wheel C
\end{definition}
\begin{proof}
  This is the declared model definition of Wheel Axle Model.
\end{proof}
\begin{definition}[Supplied Atwood Figure]
  \label{def:physics:phyx-mini-0783:phyxminiproblems-problemphyxmini0783-suppliedatwoodfigure}
  \lean{PhyXMiniProblems.ProblemPhyXMini0783.SuppliedAtwoodFigure}
  \uses{def:physics:phyx-mini-0783:phyxminiproblems-problemphyxmini0783-block, def:physics:phyx-mini-0783:phyxminiproblems-problemphyxmini0783-figurelabel, def:physics:phyx-mini-0783:phyxminiproblems-problemphyxmini0783-horizontalside}
  Qualitative evidence read from the primary raster. Pixel distances are not treated as metric data, and the bitmap contains neither numerical values nor motion arrows
\end{definition}
\begin{proof}
  This is the declared model definition of Supplied Atwood Figure.
\end{proof}
\begin{definition}[Atwood Wheel Setup]
  \label{def:physics:phyx-mini-0783:phyxminiproblems-problemphyxmini0783-atwoodwheelsetup}
  \lean{PhyXMiniProblems.ProblemPhyXMini0783.AtwoodWheelSetup}
  \uses{def:physics:phyx-mini-0783:phyxminiproblems-problemphyxmini0783-massquantity, def:physics:phyx-mini-0783:phyxminiproblems-problemphyxmini0783-lengthquantity, def:physics:phyx-mini-0783:phyxminiproblems-problemphyxmini0783-accelerationmagnitudequantity, def:physics:phyx-mini-0783:phyxminiproblems-problemphyxmini0783-forcemagnitudequantity, def:physics:phyx-mini-0783:phyxminiproblems-problemphyxmini0783-momentofinertiaquantity, def:physics:phyx-mini-0783:phyxminiproblems-problemphyxmini0783-angularaccelerationmagnitudequantity, def:physics:phyx-mini-0783:phyxminiproblems-problemphyxmini0783-block, def:physics:phyx-mini-0783:phyxminiproblems-problemphyxmini0783-verticaldirection, def:physics:phyx-mini-0783:phyxminiproblems-problemphyxmini0783-rotationdirection, def:physics:phyx-mini-0783:phyxminiproblems-problemphyxmini0783-cordmodel, def:physics:phyx-mini-0783:phyxminiproblems-problemphyxmini0783-cordwheelcontactmodel, def:physics:phyx-mini-0783:phyxminiproblems-problemphyxmini0783-wheelaxlemodel, def:physics:phyx-mini-0783:phyxminiproblems-problemphyxmini0783-suppliedatwoodfigure}
  Independent physical quantities and response observables of the apparatus. The common linear acceleration, wheel angular acceleration, and two tensions are fields constrained below by mechanics; none is defined from an answer choice
\end{definition}
\begin{proof}
  This is the declared model definition of Atwood Wheel Setup.
\end{proof}
\begin{definition}[Matches Atwood Scenario]
  \label{def:physics:phyx-mini-0783:phyxminiproblems-problemphyxmini0783-matchesatwoodscenario}
  \lean{PhyXMiniProblems.ProblemPhyXMini0783.MatchesAtwoodScenario}
  \uses{def:physics:phyx-mini-0783:phyxminiproblems-problemphyxmini0783-atwoodwheelsetup}
  The ideal Atwood-machine model and component-sign convention. Positive translation is downward for the heavier-side coordinate A, upward for B, and clockwise for wheel C in the supplied front view. These are coordinate choices, not numerical acceleration conclusions
\end{definition}
\begin{proof}
  This is the declared model definition of Matches Atwood Scenario.
\end{proof}
\begin{definition}[Matches Problem Readouts]
  \label{def:physics:phyx-mini-0783:phyxminiproblems-problemphyxmini0783-matchesproblemreadouts}
  \lean{PhyXMiniProblems.ProblemPhyXMini0783.MatchesProblemReadouts}
  \uses{def:physics:phyx-mini-0783:phyxminiproblems-problemphyxmini0783-massinkilograms, def:physics:phyx-mini-0783:phyxminiproblems-problemphyxmini0783-lengthinmeters, def:physics:phyx-mini-0783:phyxminiproblems-problemphyxmini0783-momentofinertiainkilogrammeterssquared, def:physics:phyx-mini-0783:phyxminiproblems-problemphyxmini0783-atwoodwheelsetup}
  Numerical readouts supplied in the prose. The requested wheel angular acceleration and the cord tensions do not occur in this structure
\end{definition}
\begin{proof}
  This is the declared model definition of Matches Problem Readouts.
\end{proof}
\begin{definition}[Matches Supplied Atwood Figure]
  \label{def:physics:phyx-mini-0783:phyxminiproblems-problemphyxmini0783-matchessuppliedatwoodfigure}
  \lean{PhyXMiniProblems.ProblemPhyXMini0783.MatchesSuppliedAtwoodFigure}
  \uses{def:physics:phyx-mini-0783:phyxminiproblems-problemphyxmini0783-block, def:physics:phyx-mini-0783:phyxminiproblems-problemphyxmini0783-atwoodwheelsetup}
  Primary-image evidence: B is left of wheel C, A is right of it, the wheel is suspended above both attached blocks, and A is drawn larger. No numerical mechanics result is read from the image
\end{definition}
\begin{proof}
  This is the declared model definition of Matches Supplied Atwood Figure.
\end{proof}
\begin{definition}[Matches Terrestrial Gravity Calibration]
  \label{def:physics:phyx-mini-0783:phyxminiproblems-problemphyxmini0783-matchesterrestrialgravitycalibration}
  \lean{PhyXMiniProblems.ProblemPhyXMini0783.MatchesTerrestrialGravityCalibration}
  \uses{def:physics:phyx-mini-0783:phyxminiproblems-problemphyxmini0783-accelerationinmeterspersecondsquared, def:physics:phyx-mini-0783:phyxminiproblems-problemphyxmini0783-atwoodwheelsetup}
  The multiple-choice value uses the conventional introductory-mechanics calibration g = 9.8 m/s², which is implicit rather than printed in the source
\end{definition}
\begin{proof}
  This is the declared model definition of Matches Terrestrial Gravity Calibration.
\end{proof}
\begin{definition}[Has Physical Atwood Parameters]
  \label{def:physics:phyx-mini-0783:phyxminiproblems-problemphyxmini0783-hasphysicalatwoodparameters}
  \lean{PhyXMiniProblems.ProblemPhyXMini0783.HasPhysicalAtwoodParameters}
  \uses{def:physics:phyx-mini-0783:phyxminiproblems-problemphyxmini0783-massinkilograms, def:physics:phyx-mini-0783:phyxminiproblems-problemphyxmini0783-lengthinmeters, def:physics:phyx-mini-0783:phyxminiproblems-problemphyxmini0783-accelerationinmeterspersecondsquared, def:physics:phyx-mini-0783:phyxminiproblems-problemphyxmini0783-forceinnewtons, def:physics:phyx-mini-0783:phyxminiproblems-problemphyxmini0783-momentofinertiainkilogrammeterssquared, def:physics:phyx-mini-0783:phyxminiproblems-problemphyxmini0783-atwoodwheelsetup}
  Positivity and nondegeneracy of the physical input parameters
\end{definition}
\begin{proof}
  This is the declared model definition of Has Physical Atwood Parameters.
\end{proof}
\begin{definition}[Satisfies Atwood Wheel Dynamics]
  \label{def:physics:phyx-mini-0783:phyxminiproblems-problemphyxmini0783-satisfiesatwoodwheeldynamics}
  \lean{PhyXMiniProblems.ProblemPhyXMini0783.SatisfiesAtwoodWheelDynamics}
  \uses{def:physics:phyx-mini-0783:phyxminiproblems-problemphyxmini0783-massreadout, def:physics:phyx-mini-0783:phyxminiproblems-problemphyxmini0783-lengthreadout, def:physics:phyx-mini-0783:phyxminiproblems-problemphyxmini0783-accelerationreadout, def:physics:phyx-mini-0783:phyxminiproblems-problemphyxmini0783-forcereadout, def:physics:phyx-mini-0783:phyxminiproblems-problemphyxmini0783-momentofinertiareadout, def:physics:phyx-mini-0783:phyxminiproblems-problemphyxmini0783-angularaccelerationreadout, def:physics:phyx-mini-0783:phyxminiproblems-problemphyxmini0783-atwoodwheelsetup}
  Scalar-component laws for the positive directions fixed above: m\_A g - T\_A = m\_A a for descending-side block A; T\_B - m\_B g = m\_B a for ascending-side block B; (T\_A - T\_B) r = I α for wheel C; a = r α from no slipping. The two tensions remain distinct because the wheel has nonzero moment of inertia. All equations hold in any coherent choice of mechanical units
\end{definition}
\begin{proof}
  This is the declared model definition of Satisfies Atwood Wheel Dynamics.
\end{proof}
\begin{lemma}[wheel Angular Acceleration Formula]
  \label{lem:physics:phyx-mini-0783:phyxminiproblems-problemphyxmini0783-wheelangularaccelerationformula}
  \lean{PhyXMiniProblems.ProblemPhyXMini0783.wheelAngularAccelerationFormula}
  \uses{def:physics:phyx-mini-0783:phyxminiproblems-problemphyxmini0783-massreadout, def:physics:phyx-mini-0783:phyxminiproblems-problemphyxmini0783-lengthreadout, def:physics:phyx-mini-0783:phyxminiproblems-problemphyxmini0783-accelerationreadout, def:physics:phyx-mini-0783:phyxminiproblems-problemphyxmini0783-momentofinertiareadout, def:physics:phyx-mini-0783:phyxminiproblems-problemphyxmini0783-angularaccelerationreadout, def:physics:phyx-mini-0783:phyxminiproblems-problemphyxmini0783-atwoodwheelsetup, def:physics:phyx-mini-0783:phyxminiproblems-problemphyxmini0783-hasphysicalatwoodparameters, def:physics:phyx-mini-0783:phyxminiproblems-problemphyxmini0783-satisfiesatwoodwheeldynamics}
  Eliminating the two tensions and the common linear acceleration gives the standard massive-pulley Atwood formula. This is a derived lemma, not a field of the governing-law structure
\end{lemma}
\begin{proof}
  The conclusion follows from the governing relations and figure readouts named in the statement of wheel Angular Acceleration Formula.
\end{proof}
\begin{lemma}[wheel Angular Acceleration Exact SI]
  \label{lem:physics:phyx-mini-0783:phyxminiproblems-problemphyxmini0783-wheelangularaccelerationexactsi}
  \lean{PhyXMiniProblems.ProblemPhyXMini0783.wheelAngularAccelerationExactSI}
  \uses{def:physics:phyx-mini-0783:phyxminiproblems-problemphyxmini0783-angularaccelerationinradianspersecondsquared, def:physics:phyx-mini-0783:phyxminiproblems-problemphyxmini0783-atwoodwheelsetup, def:physics:phyx-mini-0783:phyxminiproblems-problemphyxmini0783-matchesproblemreadouts, def:physics:phyx-mini-0783:phyxminiproblems-problemphyxmini0783-matchesterrestrialgravitycalibration, def:physics:phyx-mini-0783:phyxminiproblems-problemphyxmini0783-hasphysicalatwoodparameters, def:physics:phyx-mini-0783:phyxminiproblems-problemphyxmini0783-satisfiesatwoodwheeldynamics}
  With the supplied SI data, the exact acceleration is 2940/383 rad/s²
\end{lemma}
\begin{proof}
  The conclusion follows from the governing relations and figure readouts named in the statement of wheel Angular Acceleration Exact SI.
\end{proof}
\begin{definition}[Answer Choice]
  \label{def:physics:phyx-mini-0783:phyxminiproblems-problemphyxmini0783-answerchoice}
  \lean{PhyXMiniProblems.ProblemPhyXMini0783.AnswerChoice}
  The four response labels printed beside the exercise
\end{definition}
\begin{proof}
  This is the declared model definition of Answer Choice.
\end{proof}
\begin{definition}[answer Choice Radians Per Second Squared]
  \label{def:physics:phyx-mini-0783:phyxminiproblems-problemphyxmini0783-answerchoiceradianspersecondsquared}
  \lean{PhyXMiniProblems.ProblemPhyXMini0783.answerChoiceRadiansPerSecondSquared}
  \uses{def:physics:phyx-mini-0783:phyxminiproblems-problemphyxmini0783-answerchoice}
  Displayed angular-acceleration readout for each answer choice
\end{definition}
\begin{proof}
  This is the declared model definition of answer Choice Radians Per Second Squared.
\end{proof}
\begin{definition}[Matches Displayed Answer]
  \label{def:physics:phyx-mini-0783:phyxminiproblems-problemphyxmini0783-matchesdisplayedanswer}
  \lean{PhyXMiniProblems.ProblemPhyXMini0783.MatchesDisplayedAnswer}
  \uses{def:physics:phyx-mini-0783:phyxminiproblems-problemphyxmini0783-angularaccelerationinradianspersecondsquared, def:physics:phyx-mini-0783:phyxminiproblems-problemphyxmini0783-atwoodwheelsetup, def:physics:phyx-mini-0783:phyxminiproblems-problemphyxmini0783-answerchoice, def:physics:phyx-mini-0783:phyxminiproblems-problemphyxmini0783-answerchoiceradianspersecondsquared}
  An exact acceleration matches a two-decimal displayed answer when the error is below half of one hundredth. This generic rounding criterion does not single out any answer choice by definition
\end{definition}
\begin{proof}
  This is the declared model definition of Matches Displayed Answer.
\end{proof}
% --- Archon declaration topology end ---
% --- Archon physics formalization source end ---
